\documentclass{report}
\usepackage{amsmath,amssymb}
\setlength{\parindent}{0mm}
\setlength{\parskip}{1em}
\begin{document}
\begin{center}
\rule{6in}{1pt} \
{\large A. Gillman \\
{\bf Inverse medium scattering using the Hierarchical Poincar\'e-Steklov method }}

Rice University \\ Computational and Applied Mathematics \\ 6100 Main Street \\ M S 134 \\ Houston \\ Texas 77005
\\
{\tt adrianna.gillman@rice.edu}\\
C. Borges\\
L.  Greengard\end{center}

The Hierarchical Poincar\'e-Steklov (HPS) method is a recently developed
spectral discretization for partial differential equations that
naturally comes with an efficient direct solver. The HPS
method has proved extremely effective for solving free space
scattering problems with localized variable coefficients. For example,
for a variable media problem where the support of the variation was
100 wavelengths in size, the HPS method required 3.6 million unknowns to
achieve 9 digits of accuracy with the construction of the solver taking
six minutes on a desktop computer. Each solve only required three seconds. In this
talk, we present the first extension of the HPS method to inverse problems.
The problem under consideration is that of reconstructing material properties
from scattering measurements. We assume data is recovered from
many plane waves impinging the medium from multiple frequencies is given. While
this problem is ill-posed and nonlinear, it is turned into a sequence of
well-posed problems via a recursive linearization technique.
This approach requires the solution to a large number of forward scattering
problems reaching well into the mid-frequency regime making the HPS method
the ideal PDE solver.


\end{document}
